\section{Álgebra}

\questao{(OM$^3$SP 2026 - 1ª Fase) Roberta tem uma loja de roupas que trabalha exclusivamente com camisas. Ela vende todas as camisas por um mesmo preço, todas custam $R\$ \ 40,00$, e vendia cerca de $20$ camisas por dia. Ela fez uma promoção e percebeu que a cada $R\$ \ 1,00$ de desconto ela vendia $10$ camisas a mais. Qual será o valor da camisa para que Roberta tenha o faturamento máximo?}
         \begin{multicols}{4}
            \begin{enumerate}[label=\alph*)]
                \item $21$
                \item $28$
                \item $36$
                \item $45$
            \end{enumerate}
        \end{multicols}

\begin{solucoes}{1}
    \solucao
Tome como $x$ como o valor do desconto. Considere $P(x)$ a função que representa o preço e $Q(x)$ a função que representa a quantidade. Assim, teremos:
        \begin{equation*}
            V(x)=P(x)\cdot Q(x)
        \end{equation*}
            \begin{equation*}
            V(x)=(40-x)(20+10x)
        \end{equation*}
          \begin{equation*}
            V(x)=-10x^2+380x+800
        \end{equation*}
        Para haver faturamento máximo, $x_v=\frac{-b}{2\cdot a}=\frac{-380}{2\cdot (-10
        )}=\frac{-380}{-20}=19$. Sendo assim, como $x=19$, o novo valor da camisa será de $40-19=21$.\\
        \textbf{ALTERNATIVA A}
\end{solucoes}

\questao{(Elaboração própria) O Professor Samuel afirmou que $p$ é um número primo, e p$\geq5$, o valor da expressão algébrica $p^2 - 1$ será sempre um múltiplo de qual dos números abaixo?}

\begin{multicols}{4}
    \begin{enumerate}[label=\alph*) ]
        \item 10
        \item 24
        \item 13
        \item 19
    \end{enumerate}
\end{multicols}
\begin{solucoes}{2}
    \solucao
    \textbf{Testando valores} temos, como a questão afirmar que a regra vale para qualquer número primo, maior ou igual a 5,  só precisamos testar o primeiro valor número que se encaixa nessa regra para eliminar as alternativas incorretas.

    Primeiro, escolhemos o menor número primo maior ou igual a 5, que no caso é o próprio 5.

    Substituindo na expressão, temos:

    $$p^2-1 = 5^2 - 1 = 25 -1 = 24$$

    Analisando as alternativas, o 24 não é múltiplo do 10, nem do 13 e nem do 19, sendo apenas múltiplo do próprio 24. Assim, já resolvemos a questão.

    Porem, para ter certeza, vamos testar para o próximo número primeiro depois do 5, usando o 7.

    Temos o seguinte resultado:

    $$7^2 - 1 = 49 -1 = 48$$

    como 48 é 24 $\times$ 2, ele também é um múltiplo de 24, confirmando a regra.

    \textbf{ALTERNATIVA B}

    \solucao

    \textbf{Demonstração algébrica}, podemos utilizar também os produtos notáveis para realizar a solução dessa questão.

    Temos que:

    $$p^2 - 1 = (p+1)(p-1)$$

    e como \textbf{p} é um número primo maior que 2, então ele é impar. Logo (p-1) e (p+1) são dois números pares consecutivos, exemplo, se p = 7, teremos 6 e 8.\\

    Sempre que multiplicamos dois pares consecutivos, um deles obrigatoriamente é múltiplo de 2 e o outro é múltiplo de 4 e a multiplicação deles garantem que o resultado seja um múltiplo de 2 $\times$ 4 = 8. Assim, garantimos a divisibilidade por 8.\\

    para garantir a divisibilidade por 3, temos que em qualquer sequencia  de três números consecutivos como (p-1), p, (p+1), obrigatoriamente um deles é múltiplo de 3. \\

    Como p é primo maior que 3, o próprio p não pode ser múltiplo de 3, assim, (p-1) ou (p+1) será múltiplo de 3.\\

    Se a multiplicação de (p-1)(p+1) é divisível por 8 e também por 3, então obrigatoriamente tem que ser divisível  por 3 $\times$ 8 = 24.

    \textbf{ALTERNATIVA B}
\end{solucoes}
